\section{Previous Work}
Research \cite{BHTB} has shown that a good rule management system for \snort{}
will improve the workflow and efficiency of the IDS significantly, and a
prototype for such a system was built in 2012 called \bringhome{} \cite{BHTB}
(henceforth referred to as \bhtb{}). This prototype has been in use at Avdeling
BKI since 2012, but is suffering from lacks in functionality and efficiency due
to being a proof-of-concept prototype. \bhtb{} has, however, proven that
there is a significant gain from having an intuitive Graphical User Interface
(GUI) for managing rules. A GUI facilitates the work of the operator, as work is
shifted from complex operations in a CLI to automatic execution of common
routines triggered by buttons.

% TODO: improve below paragraph (Eirik)
There was another similar project conducted at \GUC{} in 2012 called
\snortmanager{} \cite{SnortManager} which also deals with management of \snort{}
rules. The \snortmanager{} project files are available as open source
software\cite{snortmanager_repo}, but no usage instructions are available.
The source code has not been maintained since its initial release and is designed for a specific company in mind. Enabling,
disabling and configuring rules is also cumbersome in this system as it is
based on plaintext arguments in ``policies'' that users must construct
themselves \cite{SnortManager}. The rule distribution is also done manually from the
files generated. All this makes \snortmanager{} unsuitable for Avdeling BKI and
their routines.

Several add-ons and other \snort-related software
exist\cite{snort_additional_downloads}, and
a few of these are built to manage \snort{} rules, most notably
\pulledpork{}\cite{pulledpork} and
\oinkmaster{}\cite{oinkmaster}. The key differences
between all software mentioned in this section and \snowman{} will be further
explained in chapter \ref{chap:existing_solutions}.